(1) 地球旋转角速度的当地竖直投影为 $\Omega_z=\Omega_E\sin\varphi$。旋转系加速度包含 $-2\boldsymbol\Omega\times\mathbf v$；在忽略地球水平旋转分量和离心修正到主阶后，$\ddot x+\omega_0^2x=2\Omega_z\dot y$、$\ddot y+\omega_0^2y=-2\Omega_z\dot x$。

(2) 取 $z=x+iy$，上述两式合并为 $\ddot z+2i\Omega_z\dot z+\omega_0^2z=0$。令 $z\propto e^{-i\omega t}$，特征方程 $-\omega^2+2\Omega_z\omega+\omega_0^2=0$，其正负圆极化频率绝对值分别为 $\sqrt{\omega_0^2+\Omega_z^2}\pm\Omega_z$。

(3) 令 $z=e^{-i\Omega_zt}w$，直接代入得到 $\ddot w+(\omega_0^2+\Omega_z^2)w=0$。因此 $w$ 在一个固定方向作谐振动，而原水平振动方向乘上 $e^{-i\Omega_zt}$，相对地面以 $-\Omega_E\sin\varphi$ 进动，这是主阶傅科进动。

(4) 在北极 $\varphi=\pi/2$ 时进动角速度为 $-\Omega_E$，一天转过一周；赤道 $\varphi=0$ 时主阶为零；南半球反向。忽略项主要为 $O(\Omega_E^2)$ 的离心变化、垂直耦合及地球有限几何，需满足 $\Omega_E/\omega_0\ll1$。
